\fsection{Contexto del proyecto}
En el contexto actual de la Industria 4.0, la automatización inteligente de procesos productivos se ha convertido en un factor
clave para la competitividad de las empresas. La irrupción del comercio electrónico y la creciente demanda de personalización en
los pedidos han generado nuevas necesidades en los centros logísticos y de distribución: mayor velocidad, flexibilidad y precisión
en la manipulación de productos. En este escenario, los sistemas tradicionales de picking manual o robotizado rígido resultan
insuficientes, ya que requieren que las piezas estén ordenadas o posicionadas de forma predefinida, lo que limita su adaptabilidad
y encarece su mantenimiento.
